\part{Dynamical Scenes}

\secn{74}{74}{Finite-horizon information and commitment}

Finite trajectories are contrast domains, but a stochastic stream must satisfy joint causal factorization, not merely causal one-coordinate marginals. Online policies obey a horizon-level entropy bound and can improve the achievable cost/value frontier. Causal attacks form a restricted feasible class whose annihilation threshold must be compared with the same unconstrained model. Provenance testing requires a specified retained-record law and information available to the forger; commitment helps only when it enforces a positive coupling gap.

\secn{75}{75}{Scenes, streams, and the filtration}

\begin{thmdef}[\hypertarget{stmt:75.1}{Definition~75.1}\ {\normalfont\upshape(dynamical scene)}]

A \textbf{dynamical scene} of horizon \(n\) consists of finite slot spaces \(X_1, \dots, X_n\) and a \textbf{trajectory domain} \(D \subseteq \prod_{t \le n} X_t\): candidates are trajectories \(Z = (Z_1, \dots, Z_n)\), and \(Z_{\le t}\) denotes the prefix. A \textbf{causal view at time \(t\)} is a view (exact: a function; graded: a channel) whose value depends on \(Z\) only through \(Z_{\le t}\); an \textbf{anchored causal view at time \(t\)} is an anchored view (Definition 19.6) whose emission and context components are prefix-measurable, its context written \(e_t(Z)\). For graded streams require a \textbf{joint causal factorization} into kernels at time \(t\) depending only on \(Z_{\le t}\) and past outputs (fresh randomization is independent of future candidate coordinates). Prefix-measurability of individual marginal channels alone does not suffice: future-dependent couplings of two null marginals would otherwise leak future information. The \textbf{record stream} is the resulting view family, graded by time; per \hyperlink{stmt:29.2}{Remark~29.2}, Part IV's reduced form is recovered as the single-slot marginal, and what the scene adds is exactly the joint structure across slots that (A7) deferred.
\end{thmdef}

\begin{thmdef}[\hypertarget{stmt:75.2}{Definition~75.2}\ {\normalfont\upshape(online policy)}]

An \textbf{online policy} for an anchored stream is a family \(\pi = (\pi_t)_{t \le n}\) where \(\pi_t\) selects the time-\(t\) verdict as a function of the \textbf{context history} \((e_1(Z), \dots, e_t(Z))\), subject to per-time compatibility as in \hyperlink{stmt:19.5}{Definition~19.5}. The policy is \textbf{memoryless} if each \(\pi_t\) depends on \(e_t\) alone. The verdict views \(\tau^\pi_t\) are causal views at time \(t\); the \textbf{verdict process} is their family. Per-round cost: \(E_t(\pi) = \mu\{Z : \pi_t \text{ errs at } Z\}\); value as in \hyperlink{sec:63}{§63} with the verdict process adjoined.
\end{thmdef}

\begin{thmdef}[\hypertarget{stmt:75.3}{Definition~75.3}\ {\normalfont\upshape(causal pattern)}]

A \textbf{causal targeted pattern} on the record stream corrupts the time-\(t\) record at rate \(\varepsilon_t\) with clutter \(Q_{Z_{\le t},\, t}\)---\hspace{0pt}measurable in the prefix (and, harmlessly, in past corruption randomness). A \textbf{clairvoyant} pattern is Part VIII's: unrestricted \(Z\)-dependence. Blind patterns are prefix-measurable trivially, so the causal/clairvoyant distinction is a targeted-regime distinction.
\end{thmdef}

\begin{thmplain}[\hypertarget{stmt:75.4}{Proposition~75.4}\ {\normalfont\upshape(transport; the filtration)}]

A dynamical scene's trajectory domain with its causal views is a contrast domain with views in the sense of Part I, so the finite constructions of Parts I--IX apply to the trajectory-domain views whenever their individual hypotheses hold. Causal restrictions constrain which channels or policies are allowed and do not disappear under this change of domain. The new object is the \textbf{filtration of registered content}: for a view set \(S\),
\[
\sigma_t(S) \;=\; \sigma\big(S \cap V_{\le t}\big), \qquad \sigma_1(S) \le \sigma_2(S) \le \cdots \le \sigma_n(S) = \sigma(S),
\]
an increasing sequence in \(\mathrm{Part}(D)\) (monotonicity of \(\sigma\)), with the graded and admissible analogues defined the same way. A question is \textbf{answerable at time \(t\)} iff \(\ker(q) \le \sigma_t(S)\) (respectively the stratum-wise and ceiling-wise forms); \textbf{prediction} is answerability at \(t\) of a question measurable only in coordinates beyond \(t\)---\hspace{0pt}a special case, not a new notion. Each increment \([\sigma_t, \sigma_{t+1}]\) is an interval of the exact calculus, so the evaluations of Parts VII and IX apply to the \emph{time} axis as they did to the cover axis: \(\nu_{\mu,q}\) decomposes the stream's information additively over rounds (chain additivity), while \(\nu_\delta\) is subadditive over rounds with the defect calculus of §69---\hspace{0pt}the quasi-cocycle theorem acquires a temporal reading for free. \(\square\)
\end{thmplain}

\begin{thmrem}[\hypertarget{stmt:75.5}{Remark~75.5}\ {\normalfont\upshape(two Part VIII phenomena were dynamics in disguise)}]

Accumulation (§23.1) is native here: each time step is a fresh look, so (K3̂)'s trade-off is not an exotic corner but the stream's default condition, and \hyperlink{stmt:61.1}{Theorem~61.1}(1)'s scheduling attack is literally temporal---\hspace{0pt}anti-correlated duty cycles across slots---\hspace{0pt}as Worked example J's instruments anticipated. Part X adds no new theorem to either; it observes that their natural habitat has now been built, and that both were provable in reduced form precisely because (A7)'s marginal already carried them (Remark 29.2's point, vindicated).
\end{thmrem}

\secn{76}{76}{Online control}

The exact anchored regime, tracked-target, determination-monotone \(K\), baseline \(B\) (record-free evidence of the stream), question \(q\), full-support \(\mu\). For an online policy \(\pi\) write \(\vec\tau = (\tau^\pi_1, \dots, \tau^\pi_n)\) and
\[\begin{gathered}V(\pi) = \nu_{\mu,q}\big(\big[\gamma_K(B),\ \gamma_K(B \cup \{\vec\tau\})\big]\big), \\ E_t(\pi) = \mu\{Z : \pi_t \text{ errs at } Z\}.\end{gathered}\]

\begin{thmplain}[\hypertarget{stmt:76.1}{Theorem~76.1}\ {\normalfont\upshape(horizon-robust Fano)}]

For every online policy---\hspace{0pt}adaptivity and memory included---\hspace{0pt}with \(m\ge2\) bounding the number of possible verdicts including \(t^*\) in every round (the constantly loyal one-symbol case is trivial):
\[
V(\pi) \;\le\; S_0 \;+\; \sum_{t=1}^{n} \Big[ h\big(E_t(\pi)\big) + E_t(\pi)\log_2(m-1) \Big]
\;\le\; S_0 + n\Big[h(\bar E) + \bar E \log_2(m-1)\Big],
\]
with \(S_0\) the scene constant of \hyperlink{stmt:63.2}{Theorem~63.2} and \(\bar E\) the average per-round error. The exchange rate of \hyperlink{stmt:63.2}{Theorem~63.2} is horizon- and adaptivity-invariant: history changes what a policy can do, not the price of doing it.
\end{thmplain}

\begin{proof}

By (K1), \(\gamma_K(B \cup \{\vec\tau\}) \le \sigma(B) \vee \bigvee_t \ker(\tau^\pi_t)\), so the ceiling block is a function of \(([Z]_{\sigma(B)}, \vec\tau(Z))\) and
\[
V(\pi) \le S_0 + H\big(\vec\tau(Z) \mid [Z]_{\sigma(B)}\big) \le S_0 + \sum_t H(\tau^\pi_t) ,
\]
by the chain rule. In a tracked-target scenario the time-\(t\) error event is exactly \(\{\tau^\pi_t \ne t^*\}\), so the grouping bound gives \(H(\tau^\pi_t) \le h(E_t) + E_t \log_2(m-1)\); concavity of \(h\) gives the per-symbol form.
\end{proof}

\begin{thmplain}[\hypertarget{stmt:76.2}{Theorem~76.2}\ {\normalfont\upshape(the adaptivity premium)}]

There is a two-round scene on which adaptive policies strictly Pareto-dominate all memoryless ones. Let \(D = \{1,2\} \times \{0,1\}^2\) (candidates \((k, b_1, b_2)\), uniform), \(G\) the global flip of \((b_1, b_2)\), \(q = b_1 \oplus b_2\), ceiling \(K_G\). Round 1: feature \(f_1 = k\) and an anchored view with context \(e_1 = k\) (verdict held loyal throughout). Round 2: feature \(v_2 = b_1\), and an anchored view whose presentation is \textbf{branch-ambiguous}: its context is
\[
e_2(Z) = \begin{cases} A & \text{if } k = 1,\, b_2 = 1 \ \text{ or } \ k = 2,\, b_1 = 0,\\ B & \text{otherwise,}\end{cases}
\]
with \(o \equiv t^*\). Then, with \(B = \{f_1, v_2\}\):

\begin{enumerate}
\def\labelenumi{\arabic{enumi}.}
\tightlist
\item
  \textbf{(Memoryless frontier.)} The four memoryless policies \(\tau(e_2)\) realize exactly \((E, V) \in \{(0,0),\, (\tfrac12, \tfrac12),\, (\tfrac12, \tfrac12),\, (1, 0)\}\): the loyal policy buys nothing, either single-cell erring policy buys \(\tfrac12\) bit at cost \(\tfrac12\), and the always-err policy has kernel \(\bot\) again---\hspace{0pt}the \(\bar d\) pathology of \hyperlink{stmt:63.3}{Proposition~63.3}, in a stream.
\item
  \textbf{(Adaptive point.)} The policy ``err on \(A\) iff \(k = 1\)''---\hspace{0pt}a function of \((e_1, e_2)\)---\hspace{0pt}achieves \((E, V) = (\tfrac14, \tfrac12)\): \textbf{the same class at half the cost}, strictly dominating the memoryless frontier.
\item
  Both frontiers respect \hyperlink{stmt:76.1}{Theorem~76.1} (\(\tfrac12 \le h(\tfrac14)\)); the premium is an efficiency gain, not a rate violation.
\end{enumerate}
\end{thmplain}

\begin{proof}

Baseline: \(\sigma(B) = \ker f_1 \vee \ker v_2\) is the \((k, b_1)\)-partition; within each branch its blocks and the \(G\)-orbits chain every candidate together, so \(\gamma_{K_G}(B) = P_G \wedge \sigma(B) = \ker(k)\) and \(I(q; [Z]_{\gamma(B)}) = 0\) (parity is branch-independent). Adaptive policy: its verdict kernel is the indicator partition of \(\{k=1, b_2 = 1\}\); adjoining it makes \(\sigma\) full on branch \(1\) and leaves branch \(2\)'s \(b_1\)-pairs, so the meet with \(P_G\) yields the two branch-\(1\) orbits and one branch-\(2\) block: \(I = \tfrac12\) bit (branch-\(1\) orbits are pure parity classes, mass \(\tfrac12\) total), \(E = \mu\{k=1, b_2=1\} = \tfrac14\). Memoryless err-on-\(A\): the kernel separates \(\{k=1, b_2=1\} \cup \{k=2, b_1=0\}\); on branch \(1\) this is the \(b_2\)-partition (full \(\sigma\), orbit blocks after the meet, contribution \(\tfrac12\)); on branch \(2\) the \(A\)-cell is the \(b_1 = 0\) pair, already inside \(\ker v_2\)---\hspace{0pt}the verdict adds \emph{no distinction there}, the branch stays one block, contribution \(0\); but the cost counts every erring candidate: \(E = \tfrac14 + \tfrac14 = \tfrac12\). Err-on-\(B\) is the \(b_2\)/\(b_1\)-complement, symmetric: \((\tfrac12, \tfrac12)\). Loyal and always-err have kernel \(\bot\). All eight values confirmed by exhaustive computation.
\end{proof}

\begin{thmplain}[\hypertarget{stmt:76.3}{Corollary~76.3}\ {\normalfont\upshape(what adaptivity buys; the anomaly's price minimized)}]

The premium's mechanism is stated by the exhibit: round-1 information (\(k\)) tells the policy \emph{which distinctions round 2 can purchase}, and the round-2 context, being branch-ambiguous, cannot. A memoryless policy must spend error mass on presentations whose verdicts buy nothing; the adaptive policy spends only where the kernel gains. Hence: \textbf{adaptivity buys frontier, never rate} (Theorem 76.1)---\hspace{0pt}its value is exactly the error mass that foresight avoids wasting, here a factor of two. The reading against \hyperlink{stmt:30.2}{Proposition~30.2} completes the control picture: the loyalty anomaly persists in streams (positive value still requires positive error---\hspace{0pt}Theorem 63.2(1) transports verbatim), but the adaptive erring policy is \emph{globally more loyal} than its memoryless rival at equal value. History is the interpreter's second lever, and unlike the first, it costs nothing. \(\square\)
\end{thmplain}

\secn{77}{77}{Causal corruption}

\begin{thmplain}[\hypertarget{stmt:77.1}{Theorem~77.1}\ {\normalfont\upshape(the causal gap)}]

Fix a finite stream and two nested closed feasible classes of attacks, causal and clairvoyant, with a compact rate parameter and an attainable annihilation constraint. In item 3 the model is whole-law contamination of one record at common rate \(\varepsilon\).

\begin{enumerate}
\def\labelenumi{\arabic{enumi}.}
\tightlist
\item
  \textbf{(Order.)} The causal annihilation threshold is \(\ge\) the corresponding clairvoyant optimum: causal patterns are a subset of clairvoyant ones.
\item
  \textbf{(Equality criterion.)} They coincide iff some optimal clairvoyant attack admits the required joint causal factorization (for one record, its clutter depends only on the available prefix).
\item
  \textbf{(Strict gap, exhibited.)} Horizon \(2\); the only record is \(R_1 = b_1\) at time \(1\); the environment coordinate \(b_2\) realizes at time \(2\) and is never recorded; prior \(b_1 \sim \mathrm{unif} \otimes b_2 \sim \mathrm{Bern}(\beta)\), \(0<\beta<1\), \(\beta\ne\tfrac12\); \(q=b_1\oplus b_2\). The honest \(q\)-conditional laws of \(R_1\) are \((1-\beta, \beta)\) and \((\beta, 1-\beta)\), so \(d = 2|1-2\beta|\) and the \textbf{clairvoyant} threshold is
  \[
  \varepsilon^*_{\mathrm{clair}} = \frac{|1-2\beta|}{1+|1-2\beta|} \qquad (= \tfrac13 \text{ at } \beta = \tfrac34),
  \]
  while the \textbf{causal} threshold is \(\tfrac12\) exactly, for every \(\beta \ne \tfrac12\). Ignorance of the future forces the adversary from mixture-merging to pointwise merging.
\end{enumerate}
\end{thmplain}

\begin{proof}

1 is set inclusion. 2: if such a clutter exists it is causal and attains the clairvoyant value; conversely attainment exhibits it. 3: clairvoyant is \hyperlink{stmt:62.1}{Lemma~62.1} with the stated \(d\). Causal: the time-\(1\) clutter is a function \(Q_{b_1}\); the perceived \(q\)-conditional difference is
\[\begin{gathered}(1-\varepsilon)\big[P(\cdot \mid q{=}0) - P(\cdot \mid q{=}1)\big] + \varepsilon \sum_{b_1}\big[P(b_1 \mid q{=}0) - P(b_1 \mid q{=}1)\big] Q_{b_1} \\ {}\qquad =\; (1-2\beta)\Big[(1-\varepsilon)(\delta_0 - \delta_1) + \varepsilon (Q_0 - Q_1)\Big],\end{gathered}\]
and for \(\beta \ne \tfrac12\) the tilt factors out entirely: annihilation requires \((1-\varepsilon)(\delta_0 - \delta_1) = \varepsilon(Q_1 - Q_0)\), whose \(L_1\) norms force \(\varepsilon \ge \tfrac12\), attained by the cross-clutter \(Q_0 = \delta_1\), \(Q_1 = \delta_0\). The clairvoyant attack that achieves \(\tfrac13\) at \(\beta = \tfrac34\) uses class-average clutter, which depends on \(b_2\)---\hspace{0pt}not prefix-measurable, as item 2 requires. Both thresholds verified by LP feasibility bisection.
\end{proof}

\begin{thmplain}[\hypertarget{stmt:77.2}{Corollary~77.2}\ {\normalfont\upshape(front-loading in the exhibit)}]

In \hyperlink{stmt:77.1}{Theorem~77.1}(3), the causal premium is \(1/2-d/(2+d)>0\) for \(0<\beta<1\), \(\beta\ne1/2\). It is a demonstrated benefit of the stated information constraint, not a universal formula for every early record. General causal thresholds require the prefix-constrained feasibility program in Appendix B. \(\square\)
\end{thmplain}

\secn{78}{78}{Provenance: commitment and shared noise}

\begin{thmplain}[\hypertarget{stmt:78.1}{Theorem~78.1}\ {\normalfont\upshape(causal conditional simulation)}]

Let \(H_t\) be the information available to the forger before writing the time-\(t\) verdict. Suppose the honest policy's conditional verdict law is a known kernel \(A_t(\cdot\mid H_t)\), simulable using permitted private randomness, and the retained time-\(t\) observation is conditionally independent of that verdict given \(H_t\). Suppose also that both constructions have the same conditional transition law to the next available history. If replacement of each verdict is allowed, a causal forger on the loyal substrate reproduces the honest policy's complete retained-record law.
\end{thmplain}

\begin{proof}

Given \(H_t\), draw a verdict from \(A_t\) independently of the retained observation. Conditional independence makes the joint round law identical. The equal transition kernel gives the same next-history law; induction proves equality of the full record law. A missing input, retained shared randomness, or a replacement-rate restriction can invalidate these hypotheses. Time alone is therefore insufficient to prevent counterfeit, but the claim is confined to this simulation model.
\end{proof}

The next theorem studies a different model in which a retained feature shares fresh randomness with the honest verdict and the forger must commit before observing it.

\begin{thmdef}[\hypertarget{stmt:78.2}{Definition~78.2}\ {\normalfont\upshape(co-registration; commitment)}]

At time \(t\) let the interpreter's perception be one draw \(y_t \sim u_t(\cdot \mid Z_{\le t})\), and let the round-record contain a \textbf{co-registered pair}: a retained feature \(g(y_t)\) and the verdict \(a(y_t)\), both functions of the \emph{same} draw. The stream is \textbf{committed} (write-ahead) if any forgery's verdict-replacement at time \(t\) may depend on \((Z_{\le t}, \text{record}_{<t}, \text{own randomness})\) but not on \(y_t\) or its functions---\hspace{0pt}the verdict is written before the co-registered feature is revealed to anything that could rewrite it.
\end{thmdef}

\begin{thmplain}[\hypertarget{stmt:78.3}{Theorem~78.3}\ {\normalfont\upshape(provenance from commitment \(+\) shared noise)}]

\par\nopagebreak

\begin{enumerate}
\def\labelenumi{\arabic{enumi}.}
\tightlist
\item
  \textbf{(Commitment is necessary.)} Without it---\hspace{0pt}the forger sees \(g(y_t)\) before writing---\hspace{0pt}conditional resampling, \(\hat a \sim \Pr(a(y_t)\in\cdot\mid g(y_t),Z_{\le t},\mathrm{past})\), reproduces the round-record's joint law exactly: forgery is perfect again.
\item
  \textbf{(Positive-gap commitment separates.)} Assume the fresh draw is independent of the forger's precommitment information conditional on the candidate history and past. With commitment, any forged verdict is conditionally independent of \(y_t\) given \((Z_{\le t}, \text{past})\), so the forged round-law is a \emph{product-form} coupling of the feature and verdict margins, and the honest and forged round-laws differ by at least \(g_t=\inf_\nu\mathrm{TV}(P_t,P_t^{\mathrm{feature}}\otimes\nu)\). Separation requires \(g_t>0\); shared-origin notation alone does not imply this. On the canonical exhibit---\hspace{0pt}\(y_t\) a uniform bit, \(g = a = y_t\)---\hspace{0pt}the honest pair is the diagonal law and every committed forgery yields \(\mathrm{TV} = \tfrac12\), \emph{regardless of the forged verdict's marginal}: for any \(\nu\), \(\mathrm{TV}(\mathrm{diag}, \mathrm{unif} \otimes \nu) = \tfrac12\).
\item
  \textbf{(Sequential detection, exponential.)} In the forensic setting---\hspace{0pt}provenance tested for a known candidate history, rounds conditionally independent given it---\hspace{0pt}for two fixed hypotheses with iid round laws \(P\ne Q\), the \(n\)-round Bayes error of the optimal audit is at most \(\tfrac12 \mathrm{BC}^n\), where \(\mathrm{BC} = \sum_y \sqrt{P(y)Q(y)} < 1\) is the round Bhattacharyya coefficient: since \(\min(a,b) \le \sqrt{ab}\) pointwise and \(\mathrm{BC}\) is multiplicative over products, \(1 - \mathrm{TV}(P^{\otimes n}, Q^{\otimes n}) \le \mathrm{BC}^n\). On the bit exhibit with forged marginal \(\nu\), \(\mathrm{BC}=(\sqrt{\nu(0)}+\sqrt{\nu(1)})/2\le2^{-1/2}\), with equality only at uniform \(\nu\): one round is inconclusive (error \(\le 0.354\)), twenty rounds pin provenance to error \(\le 2^{-11}\).
\end{enumerate}
\end{thmplain}

\begin{proof}

1: the resampled pair \((g(y_t), \hat a)\) has, by construction, the joint law of \((g(y_t), a(y_t))\) given \((Z_{\le t}, \text{past})\); induct over rounds. 2: conditional independence forces the product form; on the exhibit, the honest law puts mass \(\tfrac12\) on each diagonal cell while any product \(\mathrm{unif} \otimes \nu\) splits \(\tfrac{\nu_y}{2}\) across the column of each \(y\), and the four-cell \(L_1\) sum telescopes to \(1\) for every \(\nu\). 3: the optimal test errs with probability \(\tfrac12(1 - \mathrm{TV})\) at equal priors; apply the two displayed elementary facts. All values verified numerically.
\end{proof}

\begin{thmrem}[\hypertarget{stmt:78.4}{Remark~78.4}\ {\normalfont\upshape(the role of a coupling gap)}]

The committed audit compares its honest joint law with the forger's allowed laws. On the fair-bit example commitment forces a product law and leaves a uniform positive gap; Appendix D.7 gives the exact minimax test even for adaptive precommitted forgers. A different co-registered pair may have zero gap, and restrictions on a forger's knowledge can matter independently of commitment. The verdict concerns these declared hypotheses, not provenance in every possible attack model.
\end{thmrem}

\secn{79}{79}{Worked example L: the write-ahead audit log}

The laboratory examples can be read as three separate model-based design calculations. The branch policy of \hyperlink{stmt:76.2}{Theorem~76.2} obtains half a bit at half the deterministic memoryless error cost. The single-record model of \hyperlink{stmt:77.1}{Theorem~77.1} at \(\beta=3/4\) has causal threshold \(1/2\) versus clairvoyant threshold \(1/3\). For a fresh independent fair checksum bit, an honest matching verdict, and an all-round precommitted forger, Appendix D.7 gives equal-prior minimax audit error \(2^{-21}\) after twenty rounds. A real archive would need to establish those fresh-noise and attack assumptions before using that number; the bound is not a posterior probability of forgery for an arbitrary log.

\secn{80}{80}{Recovery: the static fiber}

\begin{thmplain}[\hypertarget{stmt:80.1}{Theorem~80.1}\ {\normalfont\upshape(horizon-one recovery)}]

At horizon one, with the complete candidate state available at that time, there is no earlier record history and no unrevealed future state. The adaptive and causal-prefix distinctions of §§76--77 disappear. A committed co-registered pair can nevertheless retain a one-round honest/forged gap: on the fair-bit example its equal-prior minimax error is \(1/4\) (Appendix D.7). Recovering the perfect-counterfeit model of \hyperlink{stmt:62.3}{Theorem~62.3} additionally removes the fresh shared-noise/commitment resource or makes its conditional law simulable by the forger.
\end{thmplain}

\begin{proof}

The first claims follow from the available histories. Substitute \(n=1\) in D.7 for the positive gap. The last condition restores the simulation hypotheses of \hyperlink{stmt:78.1}{Theorem~78.1}. A constant trajectory over several rounds need not be equivalent to a one-round experiment: independent repeated observations may accumulate information.
\end{proof}

\secn{81}{81}{Established results and scope}

The finite results include the adaptive-policy exhibit, the single-record causal threshold gap, conditional counterfeit simulation, and committed-bit testing. Appendix D.7 sharpens the latter to an exact minimax risk against adaptive precommitted forgers. Horizon one removes history and future-state restrictions but need not remove a co-registration gap. Infinite horizons, partial-round attacks, and general measurable causal optimization require additional hypotheses.
